\section{Separate a color word from its byte order}
For this framebuffer, construct the numerical word
\[
  c=(R\ll16)\mathbin{\mathrm{OR}}(G\ll8)\mathbin{\mathrm{OR}}B.
\]
Load each sample as an unsigned byte. Shift the red and green values, then join the
nonoverlapping fields. A word store places that numerical color in the display buffer.
The high byte is zero; this path does not introduce an alpha channel.

A red pixel has RGB bytes \texttt{ff 00 00} and numerical display word
\texttt{0x00ff0000}. In little-endian memory, that word occupies bytes
\texttt{00 00 ff 00}. Copying the file's three RGB bytes directly into the framebuffer
therefore does not produce the required four-byte word representation. Endianness
explains how the already constructed word is stored; it does not reverse the file's
specified component order.

For pixel index \(i\), advance the raster cursor by three and the framebuffer cursor
by four. These strides remain different even though both cursors refer to the same
logical pixel. For the reverse conversion, load one word and extract bits 23--16,
15--8 and 7--0 into three consecutive output bytes. Retain the original row order.

First inspect the four words for the small fixture: \texttt{00ff0000},
\texttt{0000ff00}, \texttt{000000ff} and \texttt{00ffffff}. Then inspect their stored
bytes. This gives two observations of the same result without confusing the number
with its memory representation.
